\section{Introduction}
\label{sec:intro}

Asset-pricing research asks whether observable firm characteristics predict
future stock returns. Characteristics such as asset growth, sales growth, and
share turnover can be used to form trading signals; the return difference
between firms with high and low values is called a factor premium. These
signals inform both asset-pricing tests and quantitative investment research.
Replication asks whether a published premium can be recovered from a paper's
description. This is difficult because papers rarely specify every
implementation choice. Reasonable choices about the stock sample, timing,
sorting, and portfolio weights can change both the estimated premium and its
statistical significance. A useful replication must therefore explain not only
whether two results differ, but also why.

Existing evidence shows why this distinction matters. In original-study
samples under a standardized protocol, \citet{HouXueZhang2020} find that
65.3\% of characteristics fail to meet $|t|\geq1.96$. By contrast,
\citet{ChenZimmermann2022} find that almost all characteristics can be
replicated when each paper's own design is followed. These statistics are not
a contest between two right or wrong answers: the studies impose different
implementation conventions. Sales-growth reversal provides a concrete example.
Under the standardized protocol, HXZ report a $t$-statistic of 1.08 for the
signal Chen and Zimmermann call $\mathit{MeanRankRevGrowth}$; under a
paper-specific design, C\&Z report 3.94. The contrast illustrates why a claim
that a factor has, or has not, replicated is incomplete without a description
of the implementation.

This thesis addresses that problem. It does not seek to decide which existing
replication is the definitive one. Instead, it asks which differences can be
linked to a controlled change in implementation, which remain attributable to
data, code, or other external features, and when the reported results do not
estimate the same quantity in the first place. The framework in
Section~\ref{sec:framework} separates these cases. A comparison is not forced
into a numerical gap when the original paper and a replication estimate
different objects.

The empirical design follows a simple principle: hold the paper
interpretation fixed, then vary the implementation choices that can be
examined in a common engine. Each paper is first converted into a reviewed,
evidence-backed implementation specification (MethodSpec). The language model
assists in extracting paper evidence and producing the approved signal formula;
human review resolves ambiguity. The resulting signal is then fixed while
portfolio-construction choices are changed. This design allows a direct
comparison of configurations without allowing the model to rewrite the
strategy between experiments. More detailed implementation, validation, and
data procedures are described in Section~\ref{sec:pipeline}.

The analysis examines asset growth \citep{CooperGulenSchill2008},
sales-growth reversal \citep{LakonishokShleiferVishny1994}, and share turnover
\citep{DatarNaikRadcliffe1998}. The three cases lead to different diagnoses.
Asset growth is reproduced on the paper-aligned raw-return metric, yet a
change in configuration contributes $+0.61\%$ per month to an observed
$+0.68\%$ gap on the C\&Z side, principally through the move from value to
equal weighting. Sales-growth reversal is reproduced under the C\&Z
configuration, whereas the standardized HXZ reference has the opposite sign.
For share turnover, the original paper reports a Fama--MacBeth coefficient
while the available replications report portfolio-sort spreads. Since these
are different quantities, the paper-anchored comparison is reported as
unidentified.

The thesis makes three contributions. First, it provides a framework that
separates controlled configuration effects from external differences and from
comparisons that are not meaningful. Second, the three cases show that
agreement on sign and statistical significance need not imply similar economic
magnitudes. Third, the workflow preserves the evidence and implementation
decisions behind each comparison. For researchers, this creates a record that
can be checked and extended. For practitioners, it offers a way to assess a
published factor before applying it to a particular universe and portfolio
construction rule.

Section~\ref{sec:related} places the thesis in the replication and factor
literatures. Sections~\ref{sec:framework} and \ref{sec:pipeline} develop the
framework and workflow. Sections~\ref{sec:data}--\ref{sec:design} present the
data and design; Section~\ref{sec:results} reports the evidence; and
Section~\ref{sec:discussion} discusses the implications and limitations.
